\section{Resumen ejecutivo}\label{sec:resumen}

Optimum Investments cierra 2025 con un patrimonio neto (PN) de \PEN200 mm sobre \PEN1{,}000 mm de activos
y \PEN800 mm de pasivos. Evaluada a un año con la historia de mercado 2021--2025, esa posición inicial
arriesga mucho más de lo que rinde: el resultado esperado es de solo \PEN3.96 mm, pero la volatilidad del
patrimonio es de \PEN25.4 mm y su pérdida esperada en el 5\,\% de escenarios más adversos (CVaR al 95\,\%)
alcanza \PEN50.1 mm, un cuarto del patrimonio. La cartera además incumple el propio límite de renta variable
del banco (22\,\% de los activos contra un máximo de 20\,\%).

Este informe recomienda reestructurar activos y pasivos hacia la cartera que minimiza esa pérdida de cola con
una aversión al riesgo $\lambda = 0.25$: subir la caja en soles a \PEN200 mm, concentrar la renta fija en soles
en el bono A02 (hasta su tope individual de \PEN350 mm), duplicar la posición en el bono flotante en dólares
A05 (a \PEN240 mm), liquidar el bono fijo en dólares A04 y recortar la renta variable a su piso de \PEN50 mm;
del lado de los pasivos, llevar el préstamo corto en soles L01 a su tope de \PEN280 mm y reducir el bono fijo
L02 a \PEN40 mm. Pagando \PEN2.73 mm en costos de transacción, esta cartera baja la volatilidad del patrimonio
a \PEN7.0 mm y el CVaR al 95\,\% a \PEN7.8 mm --- una caída de 84\,\% --- a cambio de solo \PEN0.85 mm de
resultado esperado neto (de \PEN3.96 mm a \PEN3.11 mm).

El principal riesgo de esta recomendación no es de mercado sino de modelo: los cuatro escenarios de estrés del
caso deprecian el sol sin excepción, y sustituir la regla de proyección de cupones flotantes por la forward del
período cambia el resultado de la cartera óptima en el escenario de estrés más severo de \mbox{$-$\PEN18.8 mm} a
\mbox{$-$\PEN5.8 mm} (Sección~\ref{sec:sensibilidad}). Además, el préstamo corto L01 llevado a su tope incumple con
certeza (100\,\% de probabilidad) el límite de vencimientos de corto plazo una vez llegado el horizonte;
exigir ese límite también en el horizonte cuesta \PEN0.29 mm de resultado esperado neto.
